\documentclass[10pt,a4paper]{amsart}
\usepackage[margin=19mm]{geometry}
\usepackage{amsmath,amssymb,mathtools}
\usepackage[scr=rsfs,cal=euler]{mathalfa}
\usepackage{xfrac}
\usepackage[unicode,hidelinks,bookmarks=false]{hyperref}
\newcommand{\ie}{i.e.\ }
\newcommand{\cf}{cf.\ }
\newcommand{\resp}{respectively\ }
\newcommand{\Subsection}{\subsection}
\providecommand{\nobreakdash}{\nobreak\hbox{-}}
\setlength{\emergencystretch}{2em}
\multlinegap0pt
\usepackage{amssymb}
\usepackage{xcolor}
\usepackage{algorithm}\usepackage{algpseudocode}
\usepackage[shortlabels]{enumitem}
\newtheorem{assumption}{Assumption}
\newtheorem{lemma}{Lemma}
\title{VeloTree: Inferring single-cell trajectories from RNA velocity fields with varifold distances}
\author{Elodie Maignant, Tim Conrad, Christoph von Tycowicz}
\date{Source: arXiv:2604.02380v1 (CC BY 4.0)}
\begin{document}
\maketitle

\noindent\emph{Source and licence.} VeloTree: Inferring single-cell trajectories from RNA velocity fields with varifold distances. Version arXiv:2604.02380v1 (CC BY 4.0), distributed by arXiv under the Creative Commons Attribution 4.0 International licence, http://creativecommons.org/licenses/by/4.0/, as recorded in the arXiv licence metadata for this exact version and shown on the abstract page https://arxiv.org/abs/2604.02380v1. Full licence notice: \texttt{licenses/arxiv-2604.02380-CC-BY-4.0.txt}.\par\smallskip

\noindent\textit{Editorial scope.} Counted unit: the article's lemma lemma:1 (source0.tex lines 368--374). The article's complete original proof of it is delivered with it (source0.tex lines 375--377, one \texttt{proof} environment of 3 lines). No mathematics is written, repaired or re-derived: the statement and the proof are the article's own lines, in the article's own notation, and no retained line is substituted or re-flowed.  Retained uncounted as the source-local context the proof needs: lines 311--325; lines 656--656.  One declared adaptation: exactly 1 ancillary strings inside the retained spans are omitted (image-inclusion commands and, where the source repeats a label, the repeated label definitions) and no other line differs from the source; the figure environments, with their placement, captions and labels, are kept, and the package records the omitted strings in fidelity receipts bound to the affected bodies.  No retained line contains a citation, so the wrapper carries no bibliography and no bibliography entry.  The retained lines contain the article's own diagram code, which the wrapper typesets with the standard tikz packages; no image file is delivered and the package declares no asset dependency.  Editorial formatting only, in addition: an AMS \texttt{amsart} preamble that restates the article's own declarations and macros used by the retained lines, namely the environments \texttt{assumption}, \texttt{lemma} on separate counters, together with the standard packages those lines need.  The retained environments print in this document as Assumption 1, Lemma 1 in order of appearance, whereas the article prints the same items with the numbering of its own full document.  The complete native source of the article as distributed in the arXiv e-print for this exact version is bundled unchanged as \texttt{sources/197730/source0.tex}.
\begin{assumption}
\label{A1} 
The distance between two adjacent edges of the same path is increasing as they move further from their common node. Let $i, j, k$ be three nodes of $T$ such that $i < j < k$ meaning $i$ is the parent node of $j$ which is itself the parent node of $k$. Then for all $x, x' \in \gamma_{x_i, x_j}$ and $y, y'\in \gamma_{x_j, x_k}$ we have that
\begin{equation*}
    \ell(\gamma_{x, y}) \leq \ell(\gamma_{x', y'}) \Rightarrow \|x - y\| \leq \|x' - y'\|
\end{equation*}   
where $\ell(\gamma)$ denotes the arc length of the curve $\gamma$. \\
\begin{figure}[!h]  
    \centering
    \vspace{-.3cm}
     
    \vspace{-1.4cm}
    \label{fig:hyp1}
\end{figure}
\end{assumption}

\section{Proofs}\label{app:proofs}

\begin{lemma}\label{lemma:1}
    Let $i, j$ and $k$ be three adjacent nodes of $T$ such that $i < j < k$. Assume that Assumption \ref{A1} holds. Then we have
\begin{equation}
    \label{eq:L1}
    \langle \mu_{\gamma_{x_i, x_j}}, \mu_{\gamma_{x_j, x_k}}\rangle_{W^\ast} \underset{\substack{\vspace{.05cm} \\ \sigma_x, \sigma_t \to 0 \vspace{.1cm} \\ \sigma_x \asymp \sigma_t}}{=} o\left(\|\mu_{\gamma_{x_i, x_j}}\|_{W^\ast}^2 + \|\mu_{\gamma_{x_j, x_k}}\|_{W^\ast}^2\right).
\end{equation}
\end{lemma}

\begin{proof}
    See Appendix~\ref{app:proofs}.
\end{proof}

\end{document}
